% Options for packages loaded elsewhere
% Options for packages loaded elsewhere
\PassOptionsToPackage{unicode}{hyperref}
\PassOptionsToPackage{hyphens}{url}
\PassOptionsToPackage{dvipsnames,svgnames,x11names}{xcolor}
%
\documentclass[
]{article}
\usepackage{xcolor}
\usepackage[top=0.75in,left=0.75in,bottom=0.75in,right=0.75in]{geometry}
\usepackage{amsmath,amssymb}
\setcounter{secnumdepth}{5}
\usepackage{iftex}
\ifPDFTeX
  \usepackage[T1]{fontenc}
  \usepackage[utf8]{inputenc}
  \usepackage{textcomp} % provide euro and other symbols
\else % if luatex or xetex
  \usepackage{unicode-math} % this also loads fontspec
  \defaultfontfeatures{Scale=MatchLowercase}
  \defaultfontfeatures[\rmfamily]{Ligatures=TeX,Scale=1}
\fi
\usepackage{lmodern}
\ifPDFTeX\else
  % xetex/luatex font selection
\fi
% Use upquote if available, for straight quotes in verbatim environments
\IfFileExists{upquote.sty}{\usepackage{upquote}}{}
\IfFileExists{microtype.sty}{% use microtype if available
  \usepackage[]{microtype}
  \UseMicrotypeSet[protrusion]{basicmath} % disable protrusion for tt fonts
}{}
\makeatletter
\@ifundefined{KOMAClassName}{% if non-KOMA class
  \IfFileExists{parskip.sty}{%
    \usepackage{parskip}
  }{% else
    \setlength{\parindent}{0pt}
    \setlength{\parskip}{6pt plus 2pt minus 1pt}}
}{% if KOMA class
  \KOMAoptions{parskip=half}}
\makeatother
% Make \paragraph and \subparagraph free-standing
\makeatletter
\ifx\paragraph\undefined\else
  \let\oldparagraph\paragraph
  \renewcommand{\paragraph}{
    \@ifstar
      \xxxParagraphStar
      \xxxParagraphNoStar
  }
  \newcommand{\xxxParagraphStar}[1]{\oldparagraph*{#1}\mbox{}}
  \newcommand{\xxxParagraphNoStar}[1]{\oldparagraph{#1}\mbox{}}
\fi
\ifx\subparagraph\undefined\else
  \let\oldsubparagraph\subparagraph
  \renewcommand{\subparagraph}{
    \@ifstar
      \xxxSubParagraphStar
      \xxxSubParagraphNoStar
  }
  \newcommand{\xxxSubParagraphStar}[1]{\oldsubparagraph*{#1}\mbox{}}
  \newcommand{\xxxSubParagraphNoStar}[1]{\oldsubparagraph{#1}\mbox{}}
\fi
\makeatother

\usepackage{color}
\usepackage{fancyvrb}
\newcommand{\VerbBar}{|}
\newcommand{\VERB}{\Verb[commandchars=\\\{\}]}
\DefineVerbatimEnvironment{Highlighting}{Verbatim}{commandchars=\\\{\}}
% Add ',fontsize=\small' for more characters per line
\usepackage{framed}
\definecolor{shadecolor}{RGB}{241,243,245}
\newenvironment{Shaded}{\begin{snugshade}}{\end{snugshade}}
\newcommand{\AlertTok}[1]{\textcolor[rgb]{0.68,0.00,0.00}{#1}}
\newcommand{\AnnotationTok}[1]{\textcolor[rgb]{0.37,0.37,0.37}{#1}}
\newcommand{\AttributeTok}[1]{\textcolor[rgb]{0.40,0.45,0.13}{#1}}
\newcommand{\BaseNTok}[1]{\textcolor[rgb]{0.68,0.00,0.00}{#1}}
\newcommand{\BuiltInTok}[1]{\textcolor[rgb]{0.00,0.23,0.31}{#1}}
\newcommand{\CharTok}[1]{\textcolor[rgb]{0.13,0.47,0.30}{#1}}
\newcommand{\CommentTok}[1]{\textcolor[rgb]{0.37,0.37,0.37}{#1}}
\newcommand{\CommentVarTok}[1]{\textcolor[rgb]{0.37,0.37,0.37}{\textit{#1}}}
\newcommand{\ConstantTok}[1]{\textcolor[rgb]{0.56,0.35,0.01}{#1}}
\newcommand{\ControlFlowTok}[1]{\textcolor[rgb]{0.00,0.23,0.31}{\textbf{#1}}}
\newcommand{\DataTypeTok}[1]{\textcolor[rgb]{0.68,0.00,0.00}{#1}}
\newcommand{\DecValTok}[1]{\textcolor[rgb]{0.68,0.00,0.00}{#1}}
\newcommand{\DocumentationTok}[1]{\textcolor[rgb]{0.37,0.37,0.37}{\textit{#1}}}
\newcommand{\ErrorTok}[1]{\textcolor[rgb]{0.68,0.00,0.00}{#1}}
\newcommand{\ExtensionTok}[1]{\textcolor[rgb]{0.00,0.23,0.31}{#1}}
\newcommand{\FloatTok}[1]{\textcolor[rgb]{0.68,0.00,0.00}{#1}}
\newcommand{\FunctionTok}[1]{\textcolor[rgb]{0.28,0.35,0.67}{#1}}
\newcommand{\ImportTok}[1]{\textcolor[rgb]{0.00,0.46,0.62}{#1}}
\newcommand{\InformationTok}[1]{\textcolor[rgb]{0.37,0.37,0.37}{#1}}
\newcommand{\KeywordTok}[1]{\textcolor[rgb]{0.00,0.23,0.31}{\textbf{#1}}}
\newcommand{\NormalTok}[1]{\textcolor[rgb]{0.00,0.23,0.31}{#1}}
\newcommand{\OperatorTok}[1]{\textcolor[rgb]{0.37,0.37,0.37}{#1}}
\newcommand{\OtherTok}[1]{\textcolor[rgb]{0.00,0.23,0.31}{#1}}
\newcommand{\PreprocessorTok}[1]{\textcolor[rgb]{0.68,0.00,0.00}{#1}}
\newcommand{\RegionMarkerTok}[1]{\textcolor[rgb]{0.00,0.23,0.31}{#1}}
\newcommand{\SpecialCharTok}[1]{\textcolor[rgb]{0.37,0.37,0.37}{#1}}
\newcommand{\SpecialStringTok}[1]{\textcolor[rgb]{0.13,0.47,0.30}{#1}}
\newcommand{\StringTok}[1]{\textcolor[rgb]{0.13,0.47,0.30}{#1}}
\newcommand{\VariableTok}[1]{\textcolor[rgb]{0.07,0.07,0.07}{#1}}
\newcommand{\VerbatimStringTok}[1]{\textcolor[rgb]{0.13,0.47,0.30}{#1}}
\newcommand{\WarningTok}[1]{\textcolor[rgb]{0.37,0.37,0.37}{\textit{#1}}}

\usepackage{longtable,booktabs,array}
\newcounter{none} % for unnumbered tables
\usepackage{calc} % for calculating minipage widths
% Correct order of tables after \paragraph or \subparagraph
\usepackage{etoolbox}
\makeatletter
\patchcmd\longtable{\par}{\if@noskipsec\mbox{}\fi\par}{}{}
\makeatother
% Allow footnotes in longtable head/foot
\IfFileExists{footnotehyper.sty}{\usepackage{footnotehyper}}{\usepackage{footnote}}
\makesavenoteenv{longtable}
\usepackage{graphicx}
\makeatletter
\newsavebox\pandoc@box
\newcommand*\pandocbounded[1]{% scales image to fit in text height/width
  \sbox\pandoc@box{#1}%
  \Gscale@div\@tempa{\textheight}{\dimexpr\ht\pandoc@box+\dp\pandoc@box\relax}%
  \Gscale@div\@tempb{\linewidth}{\wd\pandoc@box}%
  \ifdim\@tempb\p@<\@tempa\p@\let\@tempa\@tempb\fi% select the smaller of both
  \ifdim\@tempa\p@<\p@\scalebox{\@tempa}{\usebox\pandoc@box}%
  \else\usebox{\pandoc@box}%
  \fi%
}
% Set default figure placement to htbp
\def\fps@figure{htbp}
\makeatother


% definitions for citeproc citations
\NewDocumentCommand\citeproctext{}{}
\NewDocumentCommand\citeproc{mm}{%
  \begingroup\def\citeproctext{#2}\cite{#1}\endgroup}
\makeatletter
 % allow citations to break across lines
 \let\@cite@ofmt\@firstofone
 % avoid brackets around text for \cite:
 \def\@biblabel#1{}
 \def\@cite#1#2{{#1\if@tempswa , #2\fi}}
\makeatother
\newlength{\cslhangindent}
\setlength{\cslhangindent}{1.5em}
\newlength{\csllabelwidth}
\setlength{\csllabelwidth}{3em}
\newenvironment{CSLReferences}[2] % #1 hanging-indent, #2 entry-spacing
 {\begin{list}{}{%
  \setlength{\itemindent}{0pt}
  \setlength{\leftmargin}{0pt}
  \setlength{\parsep}{0pt}
  % turn on hanging indent if param 1 is 1
  \ifodd #1
   \setlength{\leftmargin}{\cslhangindent}
   \setlength{\itemindent}{-1\cslhangindent}
  \fi
  % set entry spacing
  \setlength{\itemsep}{#2\baselineskip}}}
 {\end{list}}
\usepackage{calc}
\newcommand{\CSLBlock}[1]{\hfill\break\parbox[t]{\linewidth}{\strut\ignorespaces#1\strut}}
\newcommand{\CSLLeftMargin}[1]{\parbox[t]{\csllabelwidth}{\strut#1\strut}}
\newcommand{\CSLRightInline}[1]{\parbox[t]{\linewidth - \csllabelwidth}{\strut#1\strut}}
\newcommand{\CSLIndent}[1]{\hspace{\cslhangindent}#1}



\setlength{\emergencystretch}{3em} % prevent overfull lines

\providecommand{\tightlist}{%
  \setlength{\itemsep}{0pt}\setlength{\parskip}{0pt}}



 


\makeatletter
\@ifpackageloaded{caption}{}{\usepackage{caption}}
\AtBeginDocument{%
\ifdefined\contentsname
  \renewcommand*\contentsname{Table of contents}
\else
  \newcommand\contentsname{Table of contents}
\fi
\ifdefined\listfigurename
  \renewcommand*\listfigurename{List of Figures}
\else
  \newcommand\listfigurename{List of Figures}
\fi
\ifdefined\listtablename
  \renewcommand*\listtablename{List of Tables}
\else
  \newcommand\listtablename{List of Tables}
\fi
\ifdefined\figurename
  \renewcommand*\figurename{Figure}
\else
  \newcommand\figurename{Figure}
\fi
\ifdefined\tablename
  \renewcommand*\tablename{Table}
\else
  \newcommand\tablename{Table}
\fi
}
\@ifpackageloaded{float}{}{\usepackage{float}}
\floatstyle{ruled}
\@ifundefined{c@chapter}{\newfloat{codelisting}{h}{lop}}{\newfloat{codelisting}{h}{lop}[chapter]}
\floatname{codelisting}{Listing}
\newcommand*\listoflistings{\listof{codelisting}{List of Listings}}
\makeatother
\makeatletter
\makeatother
\makeatletter
\@ifpackageloaded{caption}{}{\usepackage{caption}}
\@ifpackageloaded{subcaption}{}{\usepackage{subcaption}}
\makeatother
\usepackage{bookmark}
\IfFileExists{xurl.sty}{\usepackage{xurl}}{} % add URL line breaks if available
\urlstyle{same}
\hypersetup{
  pdftitle={Lab Thirteen -- Statistical Inference},
  colorlinks=true,
  linkcolor={blue},
  filecolor={Maroon},
  citecolor={Blue},
  urlcolor={Blue},
  pdfcreator={LaTeX via pandoc}}


\title{Lab Thirteen -- Statistical Inference}
\author{}
\date{}
\begin{document}
\maketitle


\normalsize

\normalsize

\section{Remaining Portfolio Work!}\label{remaining-portfolio-work}

\textbf{By 5p 4/27} Email me a link to your GitHub page with all
polished labs rendered, and your review material (at least one unit). In
the email, you should include a reflection about the processes involved
in creating this portfolio and what you learned as a result. If you're
having trouble reflecting, consider Fink's Taxonomy of Significant
Learning (Google Fink taxonomy aligned self-reflection questions). In
the email, let me know whether you would be comfortable sharing your
page -- I'm hoping that together we can get a good study guide for the
final.

\section{Statistical Inference}\label{statistical-inference}

Kasdin et al. (2025) show that dopamine in the brains of young zebra
finches acts as a learning signal, increasing when they sing closer to
their adult song and decreasing when they sing further away, effectively
guiding their vocal development through trial-and-error. This suggests
that complex natural behaviors, like learning to sing, are shaped by
dopamine-driven reinforcement learning, similar to how artificial
intelligence learns. You can find the paper at this link:
\url{https://www.nature.com/articles/s41586-025-08729-1}.

Note that the measurement is rather complicated. First, the measurements
are taken on syllables of a bird song. When discussing bird songs, the
term syllable refers to a single, distinct, uninterrupted sound that is
part of a possibly complex bird song.

They initially describe their measurement as follows. They compare the
sound of the developing zebra finches to the median sound of adult zebra
finches, that act as a ``goal'' sound.

\begin{quote}
"For each syllable rendition, we calculated the proximity (Euclidean distance) to the median of the adult renditions of the syllable (age $>$ 90 days) in feature space (Fig. 2a,b)."
\end{quote}

They contextualize that measurement by computing a ``relative quality''
that takes into account their current ability (e.g., are they getting
closer or further compared to their previous 14 renditions).

\begin{quote}
"Because dopamine is hypothesized to encode the relative value of performance, we defined the 'relative quality' of a syllable as proximity to the adult version of the syllable relative to the mean of the previous 14 renditions (Methods)."
\end{quote}

They define the ``closer'' measurements as the 10\% of syllable
renditions with the highest relative quality and the ``further'' as the
lowest 10\%.

\begin{quote}
"To test for online error responses, for each day we compared the dopamine activity associated with the 10% of syllable renditions with the highest relative quality with that of the 10% of renditions with the lowest relative quality (Fig. 2c–g)."
\end{quote}

Note that the data are measures of dopamine change, they use fibre
photometry, changes in the fluorescence indicate dopamine changes in
realtime. Their specific measurement considers the percent change in
flourescence in 100-ms windows between 200 and 300 ms from the start of
singing, averaged across development.

Finally, the researchers note that they adjust the observations to
correct for the increased dopamine for singing at all (whether it is
close or far).

\begin{quote}
"As previously reported, the baseline level of dopamine in Area X is known to increase during singing; to isolate the effect of syllable quality from these singing-related dopamine increases, we subtracted the mean dopamine response across all renditions of a syllable on each day from the dopamine response for each rendition."
\end{quote}

So, the random variables we are working with are: \begin{align*}
X_{\text{close}} &= \text{average \% change in fluorescence for the top 10\% syllable renditions} \\
&\quad \text{relative to distance from adult sound (corrected for singing)} \\[1ex]
X_{\text{far}} &= \text{average \% change in fluorescence for the bottom 10\% syllable renditions} \\
&\quad \text{relative to distance from adult sound (corrected for singing)}
\end{align*}

That is, the unit of measure is the syllable as measured on six birds,
and the measurement is somewhat complicated.

\subsection{Part a}\label{part-a}

Using the \texttt{pwr} package for \texttt{R} (Champely 2020), conduct a
power analysis. How many observations would the researchers need to
detect a moderate-to-large effect (\(d=0.65\)) when using
\(\alpha=0.05\) and default power (0.80) for a two-sided one sample
\(t\) test.

\textbf{Solution:}

\normalsize

\begin{Shaded}
\begin{Highlighting}[numbers=left,,]
\NormalTok{d }\OtherTok{=} \FloatTok{0.65}
\NormalTok{a }\OtherTok{=} \FloatTok{0.05}
\NormalTok{default\_power }\OtherTok{=} \FloatTok{0.80}

\NormalTok{n }\OtherTok{=} \FunctionTok{pwr.t.test}\NormalTok{(}\AttributeTok{d =}\NormalTok{ d,}
              \AttributeTok{sig.level =}\NormalTok{ a,}
              \AttributeTok{power =}\NormalTok{ default\_power,}
              \AttributeTok{type =} \StringTok{"one.sample"}\NormalTok{, }
              \AttributeTok{alternative =} \StringTok{"two.sided"}\NormalTok{)}
\FunctionTok{print}\NormalTok{(}\StringTok{"The required sample size for a two{-}sided one sample t test is:"}\NormalTok{)}
\end{Highlighting}
\end{Shaded}

\begin{verbatim}
[1] "The required sample size for a two-sided one sample t test is:"
\end{verbatim}

\begin{Shaded}
\begin{Highlighting}[numbers=left,,]
\NormalTok{(n)}
\end{Highlighting}
\end{Shaded}

\begin{verbatim}

     One-sample t test power calculation 

              n = 20.58039
              d = 0.65
      sig.level = 0.05
          power = 0.8
    alternative = two.sided
\end{verbatim}

\normalsize

\subsection{Part b}\label{part-b}

Click the link to go to the paper. Find the source data for Figure 2.
Download the Excel file. Describe what you needed to do to collect the
data for Figure 2(g). Note that you only need the \texttt{closer\_vals}
and \texttt{further\_vals}. Ensure to \texttt{mutate()} the data to get
a difference (e.g., \texttt{closer\_vals\ -\ further\_vals}).

\textbf{Solution:}

\subsection{Part c}\label{part-c}

Summarize the data.

\begin{itemize}
\tightlist
\item
  Summarize the further data. Do the data suggest that dopamine in the
  brains of young zebra finches decreases when they sing further away?
\item
  Summarize the closer data. Do the data suggest that dopamine in the
  brains of young zebra finches increases when they sing closer to their
  adult song?
\item
  Summarize the paired differences. Do the data suggest that there is a
  difference between dopamine in the brains of young zebra finches when
  they sing further away compared to closer to their adult song?
\item
  \textbf{Optional Challenge} Can you reproduce Figure 2(g)? Note that
  the you can use \texttt{geom\_errorbar()} to plot the range created by
  adding the mean \(\pm\) one standard deviation.
\end{itemize}

\textbf{Solution:}

\normalsize

\normalsize

\subsection{Part d}\label{part-d}

Conduct the inferences they do in the paper. Make sure to report the
results a little more comprehensively -- that is your parenthetical
should look something like: (\(t=23.99\), \(p<0.0001\); \(g=1.34\); 95\%
CI: 4.43, 4.60). Ensure to verify any conditions for use.\\
\textbf{Note}: Your numbers will vary slightly because they reported
two-sided test \(p\)-values instead of the specified one-sided tests.

\begin{itemize}
\tightlist
\item
  ``The close responses differed significantly from 0
  (\(p=1.63 \times 10^{-8}\)).'\,'\\
  \textbf{Note}: Use the appropriate one-sided test.
\item
  ``The far responses differed significantly from 0
  (\(p=5.17 \times 10^{-8}\)).'\,'\\
  \textbf{Note}: Use the appropriate one-sided test.
\item
  ``The difference between populations was significant
  (\(p=1.04 \times10^{-8}\)).'\,'\\
  \textbf{Note}: Use the two-sided test.
\end{itemize}

\textbf{Solution:}

\normalsize

\normalsize

\subsection{Part e}\label{part-e}

For each hypothesis test in the previous question, plot the null \(T\)
distribution and the observed \(t\) statistic. Ensure to use
\texttt{geom\_ribbon()} to highlight the rejection region and the
\(p\)-value. Additionally, add the resampling distribution for \(T\) to
the chart to demonstrate any differences between the null distribution
and the distribution suggested by the data.

\textbf{Solution:}

\normalsize

\normalsize

\section*{References}\label{bibliography}
\addcontentsline{toc}{section}{References}

\protect\phantomsection\label{refs}
\begin{CSLReferences}{1}{1}
\bibitem[\citeproctext]{ref-pwr}
Champely, Stephane. 2020. \emph{Pwr: {B}asic Functions for Power
Analysis}. \url{https://CRAN.R-project.org/package=pwr}.

\bibitem[\citeproctext]{ref-Kasdin25}
Kasdin, Jonathan, Alison Duffy, Nathan Nadler, et al. 2025. {``Natural
Behaviour Is Learned Through Dopamine-Mediated Reinforcement.''}
\emph{Nature}, 1--8.

\end{CSLReferences}




\end{document}
